\documentclass[12pt]{article}
\usepackage[utf8]{inputenc}
\usepackage[T1]{fontenc}
\usepackage{hyperref}
\title{Muji, Qila Language Resources for K-12 Curriculum Development}
\author{Compiled by Claude Sonnet 5.5 (claude-sonnet-5-5) for Andr\'as Kornai}
\date{}
\begin{document}
\maketitle

\section{Introduction}
Qila Muji (ISO 639-3: ymq; Glottocode qila1235) is a Loloish language of the Sino-Tibetan family. It is spoken in Qila and Wantan villages in southern Jinping County, Yunnan, China, and in Muong Gong, Lai Ch\^au Province, Vietnam. Joshua Project lists 1,600 speakers. Wikipedia's 55,000 figure covers the wider Muji cluster, not Qila Muji alone. Endangered Languages notes that it is used in daily communication but that Mandarin is the medium of education.

\section{Primary Educational Resources}
No primers, textbooks or school programmes appeared in the searches.

\section{Linguistic and Reference Materials}
\begin{itemize}
\item Wikipedia, ``Muji language'', which cites Jamin Pelkey's 2011 work on Phula variation as the main research source: \url{https://en.wikipedia.org/wiki/Muji_language}
\item Endangered Languages Project page for Qila Muji, with a bibliography. Its vitality rating is inconsistent across pages. \url{https://endangeredlanguages.com/lang/3176?hl=es}
\end{itemize}

\section{Religious and Community Resources}
\begin{itemize}
\item Joshua Project profile of Muji, Qila in China (alternate names Doka, Mujitsu, Phutsu): \url{https://joshuaproject.net/people_groups/20626/CH}
\end{itemize}

\section{Digital and Audio Resources}
None found in the searches.

\section{Curriculum Development Context}
Documentation is mostly scholarly and sparse. Education is in Mandarin, and no teaching materials were found. A native-language curriculum would depend on Pelkey's research and on new community-based work. Qila Muji has only a few villages of speakers across two countries.
\end{document}